% Regression fixture only. This text is not thesis research content.
\documentclass{tongji-proposal-word}
\school{计算机科学与技术学院}
\major{计算机科学与技术}
\student{1234567}{测试学生}
\thesistitle{视觉条件变化下道路场景语义分割的持续学习方法研究}{版面与分页测试}
\proposaldate{2026}{11}{17}
\begin{document}
\MakeProposal
\TJProposalSection[fixedfirst]{正文分页测试}
\newcommand{\TestParagraph}{本段仅用于验证长正文的自动分页与中文行距，不构成研究结论。连续填写内容应保持小四宋体、首行缩进两字符，并在超过一页时自然换页。}
\newcount\TestCounter
\TestCounter=0
\loop\ifnum\TestCounter<30
  \TJProposalBody{\TestParagraph\TestParagraph}
  \advance\TestCounter1
\repeat
\TJProposalSection{参考文献悬挂缩进测试}
\begin{TJProposalReferences}
\TJProposalReference{[1] 版面测试条目。这里的文字仅用于核验参考文献第二行相对第一行悬挂二十一磅，重复足够长度以产生自动换行。这里的文字仅用于核验参考文献第二行相对第一行悬挂二十一磅。}
\end{TJProposalReferences}
\TJProposalSection[fixedlast]{审核意见跨页测试}
\newcommand{\LongTestOpinion}{%
  \TestCounter=0
  \loop\ifnum\TestCounter<30
    \TestParagraph\par
    \advance\TestCounter1
  \repeat}
\TJProposalReviews{\LongTestOpinion}{本段为专业审核栏填写测试，应保留签名和日期，并使表格下边线完整。}
\end{document}
